\chapter{ÉTUDE PRÉALABLE}

\section*{Introduction}
\addcontentsline{toc}{section}{Introduction (Chapitre 1)}
% Présentation rapide du chapitre + objectif.

\section{Présentation de l’organisme d’accueil}
% Fiche signalétique + secteur + mission/valeurs + organigramme.

\section{Cadre du projet}
\subsection{Problématique}
% Problématique: situation actuelle + limites + enjeu.

\subsection{Objectifs à atteindre}
\begin{itemize}
  \item Objectif général : ...
  \item Objectifs spécifiques : ...
\end{itemize}

\section{Analyse de l’existant}
\subsection{Étude de l'existant}
% Processus actuel: comment les interventions sont gérées aujourd’hui ?

\subsection{Critiques de l'existant}
% Points faibles: lenteur, absence de traçabilité, erreurs, etc.

\section{Solution proposée}
% Description synthétique de la solution.

\section{Processus de développement}
\subsection{Les méthodologies agile}
% Bref rappel.

\subsection{La méthodologie SCRUM}
% Rôles, artefacts, cérémonies.

\subsection{Langages de modélisation}
% UML: cas d’utilisation, classes, séquences, etc.

\section*{Conclusion}
\addcontentsline{toc}{section}{Conclusion (Chapitre 1)}
% Résumé + transition vers chapitre suivant.